% SGA 3, Expose VI_B, Loop 1 English reconstruction.
% Scope: Section 7 heading and Proposition 7.1 in full, including
% its complete proof and source editor's notes 68--69.
% Authority: Polo--Gille Exp6B-13oct24.pdf, component-local pages
% 50--51 (combined-reader pages 391--392; re-edition headers
% 380--381; original-margin printed page 363 at the seam).
% Corollary 7.1.1 begins on local page 51 and is excluded.
% OCR and Reinhold English Markdown are comparison witnesses only.
% This unit contains no diagrams.

\section*{7. 生成的子群; 换位子群}
\phantomsection\label{sec:sga3-VIB-7}%
\addcontentsline{toc}{section}{7. 生成的子群; 换位子群}

本节始终用 \(k\) 表示一个固定的域.

\medskip
\phantomsection\label{prop:sga3-VIB-7.1}%
\noindent\textbf{命题 7.1.---}
\emph{设 \(G\) 为 \(k\)-群,
\((X_i)_{i\in I}\) 为一族几何约化的
\(k\)-概形,}%
\footnote{\label{note:sga3-VIB-68}编者注: 在此处及下文中,
将不常用的术语 ``可分'' 替换为
惯用术语 ``几何约化''; 见 EGA~IV\(_2\),
4.6.2.}
\emph{并且对每个 \(i\in I\), 设 \(f_i:X_i\to G\) 为
\(k\)-态射.}
\begin{enumerate}
\item[\textup{\SGAInlineRefTarget{prop:sga3-VIB-7.1:item:i}(i)}]\phantomsection\label{concl:sga3-VIB-7.1-i}
  \emph{存在 \(G\) 的一个最小的闭 \(k\)-子群概形,
  使得每个 \(f_i\) 都经由它分解. 将其记为
  \[
    \Gamma_G\bigl((f_i)_{i\in I}\bigr).
  \]
  它是几何约化的 \(k\)-概形, 因而当 \(G\) 在 \(k\) 上局部有限型时,
  它平滑
  (\hyperref[cor:sga3-VIB-1.3.1]{1.3.1}).}
\item[\textup{\SGAInlineRefTarget{prop:sga3-VIB-7.1:item:ii}(ii)}]\phantomsection\label{concl:sga3-VIB-7.1-ii}
  \emph{令
  \[
    X=\coprod_{i\in I}X_i,
  \]
  并设 \(f:X\to G\) 为满足对每个 \(i\in I\),
  在 \(X_i\) 上的限制为 \(f_i\) 的态射. 令
  \[
    X^1=X\amalg X,
  \]
  并设 \(f^1:X^1\to G\) 为在 \(X\) 的两个副本上的限制
  分别为 \(f\) 和 \(c\circ f\) 的态射.
  对每个 \(n>1\), 令
  \[
    X^n=X^1\times_kX^{n-1},
    \qquad
    f^n=\mu\circ(f^1\times_k f^{n-1}):X^n\longrightarrow G.
  \]
  则 \(\Gamma_G((f_i)_{i\in I})\) 是 \(G\) 的约化子概形,
  其底空间为
  \[
    \bigcup_{n\geq1}f^n(X^n).
  \]
  的闭包.}
\item[\textup{\SGAInlineRefTarget{prop:sga3-VIB-7.1:item:iii}(iii)}]\phantomsection\label{concl:sga3-VIB-7.1-iii}
  \emph{对每个 \(k\)-概形 \(S\),
  \(\Gamma_G((f_i)_{i\in I})_S\) 是 \(G_S\) 的最小闭子群
  概形, 使得每个
  \[
    f_{i,S}:(X_i)_S\longrightarrow G_S
  \]
  都经由它分解.}
\item[\textup{(iii\(^{\prime}\))}]
  \phantomsection\label{concl:sga3-VIB-7.1-iiiprime}
  \emph{此外, \(\Gamma_G((f_i)_{i\in I})_S\) 是
  \(G_S\) 的最小的、未必闭的子群概形,
  使得每个 \(f_{i,S}\) 都经由它分解.}%
  \footnote{\label{note:sga3-VIB-69}编者注: 原文在额外假设
  \(G\) 在 \(k\) 上局部有限型时陈述
  \textup{(iii\(^{\prime}\))}; 根据
  \hyperref[sga3:VIA:0.5.2]{VI\(_A\), 0.5.2}, 此假设可省略.}
\end{enumerate}

首先注意, 为证明 \textup{\SGARefLink{prop:sga3-VIB-7.1:item:i}{(i)}}, \textup{\SGARefLink{prop:sga3-VIB-7.1:item:iii}{(iii)}} 以及
\textup{(iii\(^{\prime}\))}, 按 \textup{\SGARefLink{concl:sga3-VIB-7.1-ii}{(ii)}} 定义 \(X\) 和 \(f\) 后,
即可化归为 \(I\) 只有一个元素的情形.

设 \(H\) 为 \(G\) 的约化子概形, 其底集
为
\[
  \overline{\bigcup_{n\geq1}f^n(X^n)}.
\]
态射族
\[
  f^n:X^n\longrightarrow H
\]
是概形支配的 (参见 EGA~IV\(_3\), 11.10.4), 因而
\(G\) 的每个使所有 \(f^n\) 经由它分解的闭子概形
都包含 \(H\). 此外, 由\emph{同上},
11.10.7, \(H\) 几何约化. 因此为证明
\textup{\SGARefLink{prop:sga3-VIB-7.1:item:i}{(i)}} 和 \textup{\SGARefLink{prop:sga3-VIB-7.1:item:ii}{(ii)}}, 只需证明
\(H\) 为 \(G\) 的子群概形, 即 \(c\) 在 \(H\) 上的限制
及 \(\mu\) 在
\(H\times_kH\) 上的限制均经由包含态射 \(H\to G\) 分解.

由于 \(H\) 几何约化, \(H\times_kH\) 约化,
只需在集合意义下验证
\[
  c(H)\subset H,
  \qquad
  \mu(H\times_kH)\subset H.
\]
由 EGA~IV\(_3\), 11.10.6, 所有
\[
  f^n_{(H)}(X^n\times_kH)
\]
的并在 \(H\times_kH\) 中概形稠密. 同样, 对每个
\(n\geq1\), 所有
\[
  f^m_{(X^n)}(X^n\times_kX^m),
  \qquad m\geq1,
\]
的并在 \(X^n\times_kH\) 中概形稠密.
因此只需证明
\[
  \mu\!\left(
    f^n_{(H)}
      \bigl(f^m_{(X^n)}(X^n\times_kX^m)\bigr)
  \right)\subset H
\]
及 \(c(f^n(X^n))\subset H\). 但有
\[
\begin{aligned}
  \mu\!\left(
    f^n_{(H)}
      \bigl(f^m_{(X^n)}(X^n\times_kX^m)\bigr)
  \right)
  &=
  \mu\!\left((f^n\times_k f^m)(X^n\times_kX^m)\right)\\
  &=f^{n+m}(X^{n+m})\subset H.
\end{aligned}
\]
此外, 由于 \(c(f^1(X^1))\subset f^1(X^1)\), 对每个 \(n\)
有
\[
  c(f^n(X^n))\subset f^n(X^n)\subset H.
\]
这就证明了 \textup{\SGARefLink{prop:sga3-VIB-7.1:item:i}{(i)}} 和 \textup{\SGARefLink{prop:sga3-VIB-7.1:item:ii}{(ii)}}.

现在证明 \textup{\SGARefLink{concl:sga3-VIB-7.1-iii}{(iii)}}. 设 \(G'\) 为 \(G_S\) 的闭子群
概形, 使得 \(f_S\) 经由它分解. 须证明
\(G'\) 包含 \(H_S\), 或等价地证明
\[
  H_S=H_S\times_{G_S}G'.
\]
令
\[
  H'_S=H_S\times_{G_S}G'=G'\cap H_S.
\]
每个 \(f^n_S\) 都经由 \(G'\) 和 \(H_S\) 分解,
从而经由 \(H'_S\) 分解. 由于态射族
\[
  f^n:X^n\longrightarrow H
\]
为概形支配的, EGA~IV\(_3\), 11.10.6 蕴涵
经基变换得到的态射族
\[
  f^n_S:X^n_S\longrightarrow H_S
\]
也为概形支配的. 因为每个 \(f^n_S\)
都经由 \(H'_S\) 分解, EGA~IV\(_3\), 11.10.1\textup{(c)}
给出
\[
  H'_S=H_S.
\]
这就证明了 \textup{\SGARefLink{prop:sga3-VIB-7.1:item:iii}{(iii)}}.

最后证明 \(H_S\) 为 \(G_S\) 的最小的、
未必闭的子群概形, 使得 \(f_S\)
经由它分解. 设 \(G'\) 为这样的子群概形, 并再次令
\[
  H'_S=H_S\times_{G_S}G'.
\]
须证明 \(H'_S=H_S\). 由 \textup{\SGARefLink{concl:sga3-VIB-7.1-iii}{(iii)}}, 只需
证明 \(H'_S\) 在 \(H_S\) 中闭. 为此又只需
证明 \(H_S\) 和 \(H'_S\) 的底集相同,
而为此只需对每个 \(s\in S\) 证明
\[
  H'_s
  :=H'_S\times_S\kappa(s)
  =H_s\times_{G_s}G'_s
  =H_s.
\]
由 \hyperref[sga3:VIA:0.5.2]{VI\(_A\), 0.5.2},
\(\kappa(s)\)-子群概形 \(G'_s\) 在 \(G_s\) 中闭.
所以 \(H'_s\) 在 \(H_s\) 中闭, 将上述论证
应用于 \(H'_s\), \(H_s\) 和 \(f^n_s\), 可知
\[
  H'_s=H_s.
\]
这就证明了 \textup{(iii\(^{\prime}\))}, 并完成了证明.

% 翻译统计: 源行数 200; 输出行数 201 (含本统计行); 行内/行间数学片段 120; 翻译占位 0. 底本疑点另报主线程.
